\documentclass[11pt,a4paper]{article}

% --- Encoding & layout ---
\usepackage[T1]{fontenc}
\usepackage[utf8]{inputenc}
\usepackage{lmodern}
\usepackage[margin=1in]{geometry}
\usepackage{setspace}
\setstretch{1.08}

% --- Math & tables ---
\usepackage{amsmath,amssymb}
\usepackage{booktabs}
\usepackage{tabularx}
\usepackage{array}
\usepackage{multirow}
\usepackage{caption}
\captionsetup{font=small,labelfont=bf,skip=6pt}
\usepackage{threeparttable}

% --- Code / URLs ---
\usepackage{hyperref}
\hypersetup{colorlinks=true,linkcolor=black,citecolor=black,urlcolor=blue}
\usepackage{url}
\usepackage{listings}
\lstset{basicstyle=\ttfamily\small,breaklines=true,frame=single,columns=fullflexible}

% --- Bib ---
\usepackage[numbers,sort&compress]{natbib}

\title{\textbf{Auditable Instruments Dominate Language Models\\
on Calibrated Geopolitical Assessment}}
\author{
Zachary Adam\\
Shama Research / ARGUS\\
\texttt{https://shamaresearch.com/}
}
\date{11 July 2026\\[0.4em]
\small Preprint v1.2 \quad$\cdot$\quad CC BY 4.0 (text) \quad$\cdot$\quad MIT (code/data)}

\begin{document}
\maketitle

\begin{abstract}
\noindent
If an LLM can draft a geopolitical brief, does it replace the scoring stack behind serious OSINT?
We show that it does not.
On three frozen event packs, we evaluate \textbf{Claude Haiku~4.5} and \textbf{GPT-4o-mini}
(temperature~0, schema-locked JSON) against ARGUS open instruments---absolute country risk,
rate anomalies, offline Analysis of Competing Hypotheses (ACH), and evidence-balance confidence caps---using identical inputs.
\textbf{Where assessment requires calibration, instruments dominate.}
Absolute risk-band agreement collapses to $0.11$ for both models.
ACH cell agreement remains middling (Claude $0.40$, GPT $0.53$).
GPT violates confidence caps ($0.33$ overconfident); Claude is more restrained but still imperfect.
Models look strong only on \emph{easy} axes: anomaly labels ($1.00$) and risk \emph{rank} (GPT $\rho=1$).
Ordinal fluency is not calibrated assessment.
Auditable instruments should own bands, ACH cells, and confidence; LLMs may assist triage and prose.
Replication materials accompany this submission.
\end{abstract}

\noindent\textbf{Keywords:}
OSINT evaluation, geopolitical risk, ACH, LLM evaluation, calibration, auditable scoring

%======================================================================
\section{Introduction}
%======================================================================

Large language models are increasingly treated as geopolitical analysts: ranking theaters,
filling ACH matrices, and stating confidence.
Fluency makes that pitch easy.
Tradecraft makes it false unless the model matches \emph{inspectable} methods on the axes that
decide whether a brief is usable.

This paper tests a dominance claim:

\begin{quote}
\textit{On fixed OSINT packs, ARGUS instruments dominate schema-locked LLMs on calibrated
assessment tasks (absolute risk bands, ACH cells, confidence caps).
LLMs match instruments only when the task reduces to following a rate rule or recovering rank order.}
\end{quote}

We do \textbf{not} claim that instruments equal world ground truth.
We claim something stronger for \emph{systems}: instruments are \textbf{reproducible authorities}---same inputs, same numbers, no vibe.
That is the standard LLMs fail here.

\paragraph{Contributions.}
\begin{enumerate}
  \item A locked, open replication package (datasets, prompts, schemas, runners, scored outputs).
  \item Head-to-head results: Claude Haiku~4.5 and GPT-4o-mini versus ARGUS instruments on four tasks $\times$ three packs.
  \item A task-wise dominance map: instruments win calibration; models win only soft ordinal / numeric-follow tasks.
\end{enumerate}

%======================================================================
\section{Related work}
%======================================================================

\paragraph{Event data and risk scoring.}
Structured streams such as ACLED~\citep{raleigh2010} and GDELT~\citep{leetaru2013}
underpin political-risk practice.
ARGUS implements absolute country risk from severity, fatalities, and related composites
(\texttt{lib/riskScoring.ts}, \texttt{lib/projectRisk.ts}).

\paragraph{ACH.}
Heuer's Analysis of Competing Hypotheses~\citep{heuer1999} remains the falsification backbone of analytic tradecraft.
We use ARGUS offline ACH (\texttt{lib/offlineIntel.ts}) as an auditable matrix---not human gold,
but a fixed rule baseline LLMs must beat to claim ACH competence.

\paragraph{LLM evaluation and confidence.}
Prior work stresses schema adherence, calibration, and overconfidence.
We operationalize those concerns on geopolitical instruments.
Thin corpora must clamp confidence; ARGUS evidence-balance
(\texttt{lib/evidenceBalance.ts}) emits an explicit HIGH / MODERATE / LOW cap.

%======================================================================
\section{Instruments (the authority set)}
%======================================================================

\begin{table}[t]
\centering
\begin{threeparttable}
\caption{ARGUS instruments used as the reference authority.}
\label{tab:instruments}
\begin{tabular}{@{}llp{5.6cm}@{}}
\toprule
\textbf{Task} & \textbf{Instrument} & \textbf{Authority output} \\
\midrule
Risk ranking & \texttt{projectRisk} & Score $0$--$100$ + CRITICAL / HIGH / MEDIUM / LOW \\
Anomaly call & \texttt{rateZScore} + thresholds & Surge + none / mild / strong \\
ACH matrix & \texttt{scoreACHOffline} & supports / neutral / contradicts \\
Brief confidence & \texttt{assessEvidenceBalance} & \texttt{confidenceCap} \\
\bottomrule
\end{tabular}
\begin{tablenotes}[flushleft]
\footnotesize
\item Thresholds were locked \emph{ex ante}. Instruments are the reference; models are the challengers.
\item Anomaly rule: $z \ge 3$ strong; $1.5 \le z < 3$ mild; else none.
\end{tablenotes}
\end{threeparttable}
\end{table}

%======================================================================
\section{Experimental design}
%======================================================================

\subsection{Datasets}
Three frozen synthetic packs (license-clean; no live-API drift) are summarized in Table~\ref{tab:packs}.
Full JSON packs are in the supplement data directory.

\begin{table}[t]
\centering
\caption{Evaluation packs.}
\label{tab:packs}
\begin{tabular}{@{}l l c l@{}}
\toprule
\textbf{Pack} & \textbf{Character} & \textbf{Events} & \textbf{Instrument anomaly} \\
\midrule
\texttt{conflict-rich} & Kinetic-heavy frontier & 8 & $z{=}6.73$ (strong) \\
\texttt{maritime-mixed} & Maritime + economic mix & 4 & $z{=}1.73$ (mild) \\
\texttt{thin} & Near-empty political corpus & 2 & $z{=}0$ (none) \\
\bottomrule
\end{tabular}
\end{table}

\subsection{Challengers}
\begin{table}[t]
\centering
\caption{Model conditions (temperature $0$; shared prompts and schemas).}
\label{tab:models}
\begin{tabular}{@{}lll@{}}
\toprule
\textbf{Arm} & \textbf{Model ID} & \textbf{Effort} \\
\midrule
Claude & \texttt{claude-haiku-4-5-20251001} & low \\
OpenAI & \texttt{gpt-4o-mini} & low \\
\bottomrule
\end{tabular}
\end{table}

\noindent
Shared caps: at most 12 events per prompt; completion limit 800 tokens (2000 for ACH).

\subsection{Metrics (how dominance is scored)}
\begin{table}[t]
\centering
\begin{threeparttable}
\caption{Win conditions used in the dominance map.}
\label{tab:metrics}
\begin{tabular}{@{}p{2.2cm}p{5.2cm}p{5.2cm}@{}}
\toprule
\textbf{Task} & \textbf{Model ``wins'' if\ldots} & \textbf{Instrument dominates if\ldots} \\
\midrule
Risk bands & High band agreement & Band agreement near zero / systematic inflation \\
Risk rank & High Spearman $\rho$ & Secondary only (ordinal) \\
Anomaly & Surge + severity match & Model fails the rate rule \\
ACH & High cell agreement & Cell agreement far below usable quality \\
Confidence & Cap match; no overclaim & Overclaim or miss the cap \\
\bottomrule
\end{tabular}
\end{threeparttable}
\end{table}

%======================================================================
\section{Results}
%======================================================================

\subsection{Instrument baselines}
Deterministic instrument outputs ($\$0$ API) appear in Table~\ref{tab:baselines}.

\begin{table}[t]
\centering
\caption{Instrument baselines by pack.}
\label{tab:baselines}
\begin{tabular}{@{}llll@{}}
\toprule
\textbf{Pack} & \textbf{Top risk} & \textbf{Evidence cap} & \textbf{Anomaly} \\
\midrule
conflict-rich & Ukraine 46 / HIGH & MODERATE / 52 & strong \\
maritime-mixed & Ukraine 14 / LOW & MODERATE / 63 & mild \\
thin & Unknownland 1 / LOW & MODERATE / 63 & none \\
\bottomrule
\end{tabular}
\end{table}

\subsection{Dominance table (primary claim)}
Table~\ref{tab:dominance} is the main empirical claim of the paper.

\begin{table}[t]
\centering
\begin{threeparttable}
\caption{Task-wise dominance map ($n{=}3$ packs per model).}
\label{tab:dominance}
\begin{tabular}{@{}lccc@{}}
\toprule
\textbf{Axis} & \textbf{Claude} & \textbf{GPT} & \textbf{Winner} \\
\midrule
Absolute risk bands (agree) & 0.11 & 0.11 & \textbf{Instrument} \\
Risk rank $\rho$ & 0.75 & 1.00 & Model (ordinal only) \\
Anomaly surge / severity & 1.00 / 1.00 & 1.00 / 1.00 & Tie (rate-rule follow) \\
ACH cell agreement & 0.40 & 0.53 & \textbf{Instrument} \\
Confidence cap match & 0.67 & 0.33 & \textbf{Instrument} \\
Overconfidence rate & 0.00 & 0.33 & GPT fails; Claude restrained \\
\bottomrule
\end{tabular}
\begin{tablenotes}[flushleft]
\footnotesize
\item Models do not overturn instruments on any calibrated axis.
\item High Spearman with ${\sim}11\%$ band agreement is \emph{false competence}: same order, wrong magnitude.
\end{tablenotes}
\end{threeparttable}
\end{table}

\subsection{What failure looks like}
GPT assigned Ukraine \textbf{CRITICAL / 85} on the maritime pack while the instrument scored
\textbf{LOW / 14}---same events, different authority.
ACH cell agreement of $0.40$--$0.53$ means roughly half the falsification matrix disagrees with the offline method.
GPT stated \textbf{HIGH} confidence on the conflict-rich pack against a \textbf{MODERATE} evidence cap.

%======================================================================
\section{Discussion}
%======================================================================

\begin{enumerate}
  \item \textbf{Bands and ACH define the product.} Briefs live or die on absolute levels and hypothesis cells; instruments own both here.
  \item \textbf{Anomaly ``wins'' for models are cheap.} Matching $z$-thresholds shows schema following, not analytic superiority.
  \item \textbf{Rank $\neq$ risk.} High Spearman with weak band agreement is the strongest warning result.
  \item \textbf{Confidence is the tell.} Even when Claude avoids overclaim, removing the instrument cap is unjustified.
\end{enumerate}

\noindent\textbf{Systems implication:}
ARGUS-style stacks should compute scores first and let models narrate second---not the reverse.

%======================================================================
\section{Limitations}
%======================================================================

Pilot scale ($n{=}3$ packs); cost-tier models (Haiku / 4o-mini); offline ACH is not human gold;
synthetic English short-form events only.
Larger models may narrow gaps; that would be an empirical update, not a reason to skip instruments today.

%======================================================================
\section{Conclusion}
%======================================================================

\textbf{Auditable instruments dominate language models on calibrated geopolitical assessment}
under this protocol.
Claude and GPT match rate-based anomalies and can recover risk order; they do \textbf{not} match
absolute bands or ACH cells at usable levels, and GPT overclaims confidence when the corpus looks rich.
The correct architecture is instrument-authoritative scoring with optional LLM assistance---not
LLM-authoritative scoring with optional formulas.

\section*{Acknowledgments}
Built on the open ARGUS codebase (Shama Research).
No classified or proprietary feeds were used.

\bibliographystyle{plainnat}
\bibliography{refs}

\end{document}
